\documentclass[a4paper]{article}
\usepackage[pdftex]{graphicx}
\usepackage[utf8]{inputenc}
\usepackage{enumerate}
\usepackage{amssymb}
\usepackage{href-ul}
\hypersetup{
	colorlinks=true,
	linkcolor=blue,
	urlcolor=blue}
\usepackage{geometry}
\geometry{a4paper, top=15mm, left=15mm, right=15mm, bottom=15mm,
	headsep=10mm, footskip=12mm}

\begin{document}
	{\bf magnetism}
	\begin{enumerate}[1.]
		
		%Task 1
		{\bf \item Current-carrying conductors in the magnetic field}\\
		Fill in the blanks correctly:\\
		\renewcommand{\baselinestretch}{2}\normalsize
		The so-called \rule{70 pt}{0.5 pt} acts on a current-carrying conductor in the magnetic field.
		Their direction can be changed by changing the direction of the \rule{70 pt}{0.5 pt}or the \rule{70 pt}{0.5 pt}. No effect on the current conductor in the magnetic field is observed if conductors and magnetic field lines \rule{140 pt}{0.5 pt}. The deflection of the current-carrying conductor in the magnetic field is greater, depending on:
		
		
		\begin{itemize}
			\item \rule{140 pt}{0.5 pt}
			\item \rule{140 pt}{0.5 pt}
			\item \rule{140 pt}{0.5 pt}
		\end{itemize}
		
		\renewcommand{\baselinestretch}{1}\normalsize
		
		%Exercise 2
		{\bf \item Aluminum rod on iron rails}\\
		\begin{minipage}{0.45\textwidth}
			An aluminum rod is placed on two metal rails so that it can roll freely. The arrangement is between the pole pieces of a horseshoe magnet as shown in the sketch.
		\end{minipage}
		\hfill
		\begin{minipage}{0.5\textwidth}
			\includegraphics[width=7 cm]{magnschi132}
		\end{minipage}
		
		\begin{enumerate}[a)]
			\item Briefly explain what is observed when the circuit is closed (reason).
			\item What does a reversal of the current direction cause?
		\end{enumerate}
		
		%Task 3
		{\bf \item electric motor}
		
		\includegraphics[width=13 cm]{mgnlsch132}
		
		\begin{enumerate}[a)]
			\item Draw the force arrows in stations I and II and mark the direction of rotation of the conductor loop.
			\item Enter the polarity of the conductor loop in situation III so that the direction of rotation is maintained.
			\item Explain how the electric motor works, using and explaining the terms rotor, stator, dead center and pole inverter.
		\end{enumerate}
		
		\begin{minipage}{0.55\textwidth}
			{\bf \item A conductor loop should move in the specified direction.} How should the power source be polarized? Explain your reasoning using the three finger rule.
		\end{minipage}
		\hfill
		\begin{minipage}{0.4\textwidth}
			\includegraphics[width=5 cm]{leischau214}
		\end{minipage}
		
		{\bf \item What is meant by the Lorentz force?}
		
		{\bf \item The force effect on free charge carriers caused by a magnetic field} can be shown, for example, with the help of a filament beam tube. Describe exactly how a thread jet tube works using a labeled sketch.
		
		\begin{minipage}{0.3\textwidth}
			{\bf \item An electron enters a spatially limited, homogeneous magnetic field.} Draw the further path as well as some of the forces acting on the electron as an example. Explain your solution.
		\end{minipage}
		\hfill
		\begin{minipage}{0.65\textwidth}
			\hspace{30 pt}
			\includegraphics[width=7 cm]{elo214}
		\end{minipage}
		
		
	\end{enumerate}
	\fbox{
		\begin{minipage}{0.5\textwidth}
			Additional worksheets are available at
			
			\href{http://www.okuyakl.com}{www.okuyakl.com}
		\end{minipage}
		\hfill
		\begin{minipage}{0.4\textwidth}
			\includegraphics[width=1.5 cm]{../../viecher/zwe03}
			\includegraphics[width=2 cm]{../../viecher/afanticon1}
			
	\end{minipage}}
	
\end{document}